\documentclass[10pt,a4paper]{amsart}
\usepackage{amsmath,amssymb,amsthm,graphicx}
% For commutative diagrams we use Xy-pic
\usepackage[all]{xy}

% We use 2cell for 2-commutative diagrams.
\xyoption{2cell}
% Theorem environments.
%
\theoremstyle{plain}
\newtheorem{theorem}[subsection]{Theorem}
\newtheorem{proposition}[subsection]{Proposition}
\newtheorem{lemma}[subsection]{Lemma}

\theoremstyle{definition}
\newtheorem{definition}[subsection]{Definition}
\newtheorem{example}[subsection]{Example}
\newtheorem{exercise}[subsection]{Exercise}
\newtheorem{situation}[subsection]{Situation}

\theoremstyle{remark}
\newtheorem{remark}[subsection]{Remark}
\newtheorem{remarks}[subsection]{Remarks}

\numberwithin{equation}{subsection}

% Macros
%
\def\lim{\mathop{\mathrm{lim}}\nolimits}
\def\colim{\mathop{\mathrm{colim}}\nolimits}
\def\Spec{\mathop{\mathrm{Spec}}}
\def\Hom{\mathop{\mathrm{Hom}}\nolimits}
\def\Ext{\mathop{\mathrm{Ext}}\nolimits}
\def\SheafHom{\mathop{\mathcal{H}\!\mathit{om}}\nolimits}
\def\SheafExt{\mathop{\mathcal{E}\!\mathit{xt}}\nolimits}
\def\Sch{\mathit{Sch}}
\def\Mor{\mathop{\mathrm{Mor}}\nolimits}
\def\Ob{\mathop{\mathrm{Ob}}\nolimits}
\def\Sh{\mathop{\mathit{Sh}}\nolimits}
\def\NL{\mathop{N\!L}\nolimits}
\def\CH{\mathop{\mathrm{CH}}\nolimits}
\def\proetale{{pro\text{-}\acute{e}tale}}
\def\etale{{\acute{e}tale}}
\def\QCoh{\mathit{QCoh}}
\def\Ker{\mathop{\mathrm{Ker}}}
\def\Im{\mathop{\mathrm{Im}}}
\def\Coker{\mathop{\mathrm{Coker}}}
\def\Coim{\mathop{\mathrm{Coim}}}

% Boxtimes
%
\DeclareMathSymbol{\boxtimes}{\mathbin}{AMSa}{"02}

%
% Macros for moduli stacks/spaces
%
\def\QCohstack{\mathcal{QC}\!\mathit{oh}}
\def\Cohstack{\mathcal{C}\!\mathit{oh}}
\def\Spacesstack{\mathcal{S}\!\mathit{paces}}
\def\Quotfunctor{\mathrm{Quot}}
\def\Hilbfunctor{\mathrm{Hilb}}
\def\Curvesstack{\mathcal{C}\!\mathit{urves}}
\def\Polarizedstack{\mathcal{P}\!\mathit{olarized}}
\def\Complexesstack{\mathcal{C}\!\mathit{omplexes}}
% \Pic is the operator that assigns to X its picard group, usage \Pic(X)
% \Picardstack_{X/B} denotes the Picard stack of X over B
% \Picardfunctor_{X/B} denotes the Picard functor of X over B
\def\Pic{\mathop{\mathrm{Pic}}\nolimits}
\def\Picardstack{\mathcal{P}\!\mathit{ic}}
\def\Picardfunctor{\mathrm{Pic}}
\def\Deformationcategory{\mathcal{D}\!\mathit{ef}}
\usepackage{cleveref}
\crefname{originaldefinition}{original definition}{original definitions}
\Crefname{originaldefinition}{Original definition}{Original definitions}
\makeatletter
\newlabel{definition-cech-complex}{{03OL}{source}{}{sourcepointer-1}{}}
\newlabel{definition-cech-complex@cref}{{[originaldefinition][0][]03OL}{[1][0][]source}{}{}{}}
\makeatother
\crefname{originallemma}{original lemma}{original lemmas}
\Crefname{originallemma}{Original lemma}{Original lemmas}
\makeatletter
\newlabel{sites-cohomology-lemma-homology-complex}{{03AT}{source}{}{sourcepointer-2}{}}
\newlabel{sites-cohomology-lemma-homology-complex@cref}{{[originallemma][0][]03AT}{[1][0][]source}{}{}{}}
\makeatother
\title{The Cech chain complex of a family of morphisms is exact away from degree zero}
\author{The Stacks Project Authors}
\date{stacks-project:etale-cohomology.tex:local-snapshot-20261011; GNU FDL 1.2 or later}
\begin{document}
\maketitle
\noindent\textit{Editorial note.} All delivered mathematical text is the author's own, copied from the pinned Stacks Project chapter file etale-cohomology.tex; no statement or proof step was rewritten, reordered or removed. The displayed definition of delta in the proof carries the author's sign (-1)\textasciicircum{}p where the index i is intended; it is reproduced verbatim and not corrected.\par\medskip
\noindent\textit{Editorial note.} The preamble is the chapter's own native preamble excerpt used for this source: the Xy-pic setup at lines 14-18 and the theorem declarations with the mathematical macros at lines 154-217 of sources/stacks-etale-cohomology/preamble.tex. The quoted class line and the xr-hyper externaldocument block are not included because the assembler supplies a plain amsart wrapper; no macro was invented or redefined.\par\medskip
\noindent\textit{Editorial note.} The same-chapter reference is located at its \textbackslash{}label line inside the pinned chapter file; the cross-chapter pointer is located in the project's own tag list staged in the snapshot (sources/stacks-chow/tags.txt), at the verbatim line carrying its canonical tag.\par\medskip
\noindent\textit{Allowed standard prerequisites (editorial).}\begin{itemize}
\item Cech complex of a family and the presheaves Z\_X represented by objects. Same chapter, Definition environment labelled definition-cech-complex (source line 1477, Stacks tag 03OL); the statement refers to its notation and assumptions. Supplies the complex Z\textasciicircum{}bullet\_U of abelian presheaves; its exactness away from degree zero is the counted content.
\item Simplicial acyclicity of Z[S] <- Z[S x S] <- Z[S x S x S] <- .... Standard homological algebra; the proof exhibits an explicit nullhomotopy for the operator delta rather than citing it. Supplies the left-hand target of the reduction; the reduction itself and the homotopy are retained in the counted span.
\end{itemize}\medskip
\ifdefined\hypertarget\hypertarget{sourcepointer-1}{}\fi
\noindent\textit{Source reference.} Original definition 03OL: sources/stacks-etale-cohomology/pinned-original-source.tex, source lines 1477--1477. Same-chapter pointer named in the first line of the counted statement (line 1628) as the source of notation and assumptions; the reduction and the explicit homotopy are retained verbatim in the counted span.\par\medskip
\ifdefined\hypertarget\hypertarget{sourcepointer-2}{}\fi
\noindent\textit{Source reference.} Original lemma 03AT: sources/stacks-chow/tags.txt, source lines 4322--4322. Cross-chapter 'see for more details' pointer at proof lines 1696-1698, invoked after the homotopy is already constructed; the counted proof is self-contained.\par\medskip
\begin{lemma}
\label{lemma-exact}
Notation and assumptions as in Definition \ref{definition-cech-complex}.
The complex of abelian presheaves
\begin{align*}
\mathbf{Z}_\mathcal{U}^\bullet \quad : \quad
\bigoplus_{i_0 \in I} \mathbf{Z}_{U_{i_0}} \leftarrow
\bigoplus_{i_0, i_1 \in I} \mathbf{Z}_{U_{i_0} \times_U U_{i_1}} \leftarrow
\bigoplus_{i_0, i_1, i_2 \in I}
\mathbf{Z}_{U_{i_0} \times_U U_{i_1} \times_U U_{i_2}} \leftarrow
\ldots
\end{align*}
is exact in all degrees except $0$ in $\textit{PAb}(\mathcal{C})$.
\end{lemma}

\begin{proof}
For any $V\in \Ob(\mathcal{C})$ the complex of abelian groups
$\mathbf{Z}_\mathcal{U}^\bullet(V)$ is
$$
\begin{matrix}
\mathbf{Z}\left[
\coprod_{i_0\in I} \Mor_\mathcal{C}(V, U_{i_0})\right]
\leftarrow
\mathbf{Z}\left[
\coprod_{i_0, i_1 \in I}
\Mor_\mathcal{C}(V, U_{i_0} \times_U U_{i_1})\right]
\leftarrow \ldots = \\
\bigoplus_{\varphi : V \to U}
\left(
\mathbf{Z}\left[\coprod_{i_0 \in I} \Mor_\varphi(V, U_{i_0})\right]
\leftarrow
\mathbf{Z}\left[\coprod_{i_0, i_1\in I} \Mor_\varphi(V, U_{i_0}) \times
\Mor_\varphi(V, U_{i_1})\right]
\leftarrow
\ldots
\right)
\end{matrix}
$$
where
$$
\Mor_{\varphi}(V, U_i)
=
\{ V \to U_i \text{ such that } V \to U_i \to U \text{ equals } \varphi \}.
$$
Set $S_\varphi = \coprod_{i\in I} \Mor_\varphi(V, U_i)$, so that
$$
\mathbf{Z}_\mathcal{U}^\bullet(V)
=
\bigoplus_{\varphi : V \to U}
\left(
\mathbf{Z}[S_\varphi] \leftarrow
\mathbf{Z}[S_\varphi \times S_\varphi] \leftarrow
\mathbf{Z}[S_\varphi \times S_\varphi \times S_\varphi] \leftarrow
\ldots \right).
$$
Thus it suffices to show that for each $S = S_\varphi$, the complex
\begin{align*}
\mathbf{Z}[S] \leftarrow
\mathbf{Z}[S \times S] \leftarrow
\mathbf{Z}[S \times S \times S] \leftarrow \ldots
\end{align*}
is exact in negative degrees. To see this, we can give an explicit homotopy.
Fix $s\in S$ and define $K: n_{(s_0, \ldots, s_p)} \mapsto n_{(s, s_0,
\ldots, s_p)}.$ One easily checks that $K$ is a nullhomotopy for the operator
$$
\delta :
\eta_{(s_0, \ldots, s_p)}
\mapsto
\sum\nolimits_{i = 0}^p (-1)^p \eta_{(s_0, \ldots, \hat s_i, \ldots, s_p)}.
$$
See
Cohomology on Sites, Lemma \ref{sites-cohomology-lemma-homology-complex}
for more details.
\end{proof}

\end{document}
